\documentclass[11pt]{article}
\usepackage[margin=1in]{geometry}
\usepackage{amsmath,amssymb}
\usepackage[hidelinks]{hyperref}

\title{Literature-Implied Answer: A Monotonically Sokolov Characterization}
\author{Run \texttt{fa\_banach\_001}}
\date{June 18, 2026}

\newcommand{\wstar}{w^*}
\newcommand{\Tau}{\mathcal T}

\begin{document}
\emergencystretch=2em
\maketitle

\paragraph{Status.}
\texttt{literature\_implied\_answer}.  This packet records a later-literature
identification, not a new original proof.  The source paper is Wieslaw Kubis,
\emph{Banach spaces with projectional skeletons}, arXiv:0802.1109.  The
supporting paper is Ondrej F. K. Kalenda, \emph{Projectional skeletons and
Markushevich bases}, arXiv:1805.11901.

\paragraph{Source question.}
In the final remarks of arXiv:0802.1109, Kubis writes that for a norming
subspace \(D\subset X^*\), \(\Tau_D=\sigma(X,D)\), and asks:
\[
  \text{Find a topological property of } (X,\Tau_D)
  \text{ which says when } D \text{ generates projections in } X.
\]
This appears on page 28 of the source PDF, immediately after the discussion
that \(D\) generating projections implies \((X,\Tau_D)\) is Lindelof.

\paragraph{Supporting theorem.}
Kalenda's Theorem 4.1 in arXiv:1805.11901 states that, for a Banach space
\(X\) and a norming subspace \(D\subset X^*\), the following are equivalent,
among other formulations:
\[
  D \text{ is induced by a projectional skeleton in } X,
\]
and
\[
  D \text{ is weak}^{*}\text{-countably closed and }
  (X,\sigma(X,D)) \text{ is monotonically Sokolov}.
\]
The theorem also gives an equivalent continuous-image formulation using
monotonically Sokolov spaces.

\paragraph{Identification.}
Kubis had already proved the bridge between generating projections and induced
subspaces.  If a norming subspace \(D\) generates projections, then the
smallest linear subspace containing \(D\) and closed under weak-star closures
of countable subsets is induced by a projectional skeleton.  Hence, when
\(D\) itself is weak-star-countably closed, generation of projections is
equivalent to \(D\) being induced by a projectional skeleton.  Conversely, an
induced norming subspace generates projections by Kubis's characterization of
projectional skeletons via elementary submodels.

\paragraph{Conclusion.}
For the standard weak-star-countably-closed normalization of the problem, the
requested topological property is monotone Sokolovness:
\[
  D \text{ generates projections in } X
  \quad\Longleftrightarrow\quad
  (X,\sigma(X,D)) \text{ is monotonically Sokolov},
\]
provided \(D\subset X^*\) is norming and weak-star-countably closed.  More
precisely, the right-hand side should be read together with the
weak-star-countable closedness condition on \(D\), exactly as in Kalenda's
Theorem 4.1.

\paragraph{Scope limitations.}
This does not claim a complete answer for arbitrary norming subspaces that are
not weak-star-countably closed.  It also does not resolve the other final
questions in arXiv:0802.1109, such as the general projectional-generator
problem, the closed-subspace separable-complementation problem, or the
\(C(\omega_2+1)\) embedding problem for Plichko spaces.

\paragraph{Search evidence.}
Local cheap-index searches found no previous packet for arXiv:0802.1109 or
for the monotonically Sokolov characterization.  Web and local source searches
for exact phrases around ``topological property'', \(\Tau_D\),
\(\sigma(X,D)\), ``projectional generator'', ``generates projections'', and
``monotonically Sokolov'' found arXiv:1805.11901 as the decisive later anchor.
The relation to Kubis's exact problem is agent-identified from the two theorem
statements, rather than explicitly advertised by the supporting author.

\paragraph{References.}
\begin{itemize}
\item W. Kubis, \emph{Banach spaces with projectional skeletons}, arXiv:0802.1109.
\item O. F. K. Kalenda, \emph{Projectional skeletons and Markushevich bases},
  arXiv:1805.11901; Proc. London Math. Soc. (3) 120 (2020), no. 4, 514--586.
\end{itemize}

\end{document}
